\documentclass[a4paper,11pt]{article}

\usepackage[utf8]{inputenc}
\usepackage[T1]{fontenc}
\usepackage[romanian]{babel}
\usepackage[margin=2cm]{geometry}
\usepackage{booktabs}
\usepackage{tabularx}
\usepackage{longtable}
\usepackage{array}
\usepackage{xcolor}
\usepackage{listings}
\usepackage{tcolorbox}
\usepackage{titlesec}
\usepackage{enumitem}
\usepackage{amsmath}
\usepackage{amssymb}
\usepackage{fancyhdr}
\usepackage{hyperref}
\usepackage{multirow}
\usepackage{colortbl}

% ── Culori ──────────────────────────────────────────────────────────────────
\definecolor{primary}{RGB}{30, 80, 160}
\definecolor{accent}{RGB}{220, 50, 50}
\definecolor{lightgray}{RGB}{245, 245, 245}
\definecolor{darkgray}{RGB}{80, 80, 80}
\definecolor{codegreen}{RGB}{0, 120, 0}
\definecolor{codepurple}{RGB}{130, 0, 160}
\definecolor{tblheader}{RGB}{30, 80, 160}

% ── Stilizare titluri secțiuni ───────────────────────────────────────────────
\titleformat{\section}
  {\large\bfseries\color{primary}}
  {\thesection.}{0.5em}{}
  [\titlerule]

\titleformat{\subsection}
  {\normalsize\bfseries\color{primary!80!black}}
  {\thesubsection.}{0.5em}{}

\titleformat{\subsubsection}
  {\normalsize\bfseries\color{darkgray}}
  {}{0em}{}

% ── Cod sursă ────────────────────────────────────────────────────────────────
\lstdefinestyle{codestyle}{
  backgroundcolor=\color{lightgray},
  commentstyle=\color{codegreen},
  keywordstyle=\color{primary}\bfseries,
  stringstyle=\color{codepurple},
  basicstyle=\ttfamily\footnotesize,
  breaklines=true,
  frame=single,
  framerule=0.4pt,
  rulecolor=\color{primary!40},
  tabsize=2,
  showstringspaces=false,
  numbers=none,
}
\lstset{style=codestyle}

% ── Boxuri evidențiate ───────────────────────────────────────────────────────
\tcbuselibrary{skins, breakable}

\newtcolorbox{defbox}[1]{
  colback=primary!8,
  colframe=primary,
  fonttitle=\bfseries\small,
  title=#1,
  breakable,
  left=4pt, right=4pt, top=4pt, bottom=4pt,
}

\newtcolorbox{formulabox}{
  colback=accent!8,
  colframe=accent,
  breakable,
  left=4pt, right=4pt, top=4pt, bottom=4pt,
}

\newtcolorbox{tipbox}{
  colback=codegreen!8,
  colframe=codegreen,
  fonttitle=\bfseries\small,
  title=Atenție,
  breakable,
  left=4pt, right=4pt, top=4pt, bottom=4pt,
}

% ── Header/Footer ────────────────────────────────────────────────────────────
\pagestyle{fancy}
\fancyhf{}
\fancyhead[L]{\small\color{darkgray}\textbf{APD — Sinteză Cursuri 1–6}}
\fancyhead[R]{\small\color{darkgray}\leftmark}
\fancyfoot[C]{\small\color{darkgray}\thepage}
\renewcommand{\headrulewidth}{0.4pt}

% ── Tabele colorate ──────────────────────────────────────────────────────────
\newcommand{\thead}[1]{\cellcolor{tblheader}\color{white}\textbf{#1}}

% ── Hyperref ─────────────────────────────────────────────────────────────────
\hypersetup{
  colorlinks=true,
  linkcolor=primary,
  urlcolor=primary,
}

% ─────────────────────────────────────────────────────────────────────────────
\begin{document}

\begin{titlepage}
  \centering
  \vspace*{2cm}
  {\Huge\bfseries\color{primary} Algoritmi Paraleli și Distribuiți \par}
  \vspace{0.5cm}
  {\Large\color{darkgray} Sinteză pentru Examen Parțial \par}
  \vspace{0.3cm}
  {\large\color{darkgray} Cursurile 1–6 \par}
  \vspace{2cm}
  \hrule
  \vspace{0.5cm}
  {\normalsize Conținut: Introducere, Tipuri de Paralelism, Modele de Calcul,
  Metrici de Performanță, Arhitecturi, Proiectarea Programelor Paralele, OpenMP \par}
  \vspace{0.5cm}
  \hrule
  \vfill
\end{titlepage}

\tableofcontents
\newpage

% =============================================================================
\section{Curs 1 — Introducere în APD}
% =============================================================================

\subsection{Motivație și domenii de aplicație}

Paralelismul/distributivitatea permit:
\begin{itemize}[noitemsep]
  \item Trecerea de la \textbf{gândirea secvențială} la \textbf{gândirea paralelă}
  \item Utilizarea procesoarelor multicore moderne
  \item Accelerarea calculelor în: prognoze meteo, genetică, finanțe, inginerie spațială
\end{itemize}

\subsection{Definiții esențiale}

\begin{defbox}{Definiții fundamentale}
\begin{description}[leftmargin=3cm, style=nextline, noitemsep]
  \item[\textbf{Paralelism}] Utilizarea \emph{simultană} a mai multor resurse de calcul pentru accelerare
  \item[\textbf{Concurență}] Administrarea \emph{corectă și eficientă} a accesului la resurse de calcul
  \item[\textbf{Task/Activitate}] Un calcul care poate rula concurent cu altele (program, proces, fir)
  \item[\textbf{Atomicitate}] Proprietatea unei operații de a fi indivizibilă
  \item[\textbf{Consensus}] Acordul între două sau mai multe activități asupra unui predicat
  \item[\textbf{Consistență}] Reguli convenite despre valorile variabilelor în memoria partajată
  \item[\textbf{Multicast}] Transmiterea unui mesaj către mai mulți destinatari fără constrângeri de ordine
\end{description}
\end{defbox}

\begin{tipbox}
\textbf{Concurența $\neq$ Paralelism}: Concurența = logic active simultan; Paralelismul = fizic active simultan.
Programele paralele includ programele concurente, dar nu invers.
\end{tipbox}

\subsection{Calculator paralel vs. distribuit}

\begin{center}
\begin{tabular}{lcc}
\toprule
\thead{Caracteristică} & \thead{Paralel} & \thead{Distribuit} \\
\midrule
Scop general & Viteză & Comoditate \\
Interacțiuni & Frecvente & Rare \\
Granularitate & Mare & Mică \\
Siguranță funcționare & Presupusă adevărată & Posibil neadevărată \\
\bottomrule
\end{tabular}
\end{center}

\begin{itemize}[noitemsep]
  \item \textbf{Sistem paralel} = procesoare \emph{strâns cuplate} pe același sistem, comunică prin memorie partajată
  \item \textbf{Sistem distribuit} = sisteme \emph{autonome} în rețea, comunică prin mesaje
\end{itemize}

\subsection{Avantajele sistemelor distribuite}
\begin{enumerate}[noitemsep]
  \item Partajarea datelor și resurselor
  \item Comunicare (email, chat)
  \item Flexibilitate (distribuirea sarcinilor)
  \item Cost scăzut față de mainframe
  \item Viteză prin load balancing
  \item Distributivitate nativă (magazine, bănci)
  \item Toleranță la defecte prin redundanță
\end{enumerate}

\subsection{Limitările sistemelor distribuite}
\begin{itemize}[noitemsep]
  \item Software greu de dezvoltat
  \item Rețea: întârzieri, saturație, pierderi de pachete
  \item Securitate
  \item \textbf{Problema consensului}: nodurile trebuie să convină prin canale nesigure
\end{itemize}

\subsection{Problema celor 2 Generali}

\begin{defbox}{Problemă fără soluție}
Doi generali trebuie să atace simultan comunicând prin mesageri care pot fi uciși.
\textbf{Nu are soluție} în prezența eșecurilor arbitrare de comunicare.
\end{defbox}

\begin{formulabox}
\textbf{Teoremă cheie:} Transmiterea atomică $\equiv$ Consensus (sunt echivalente!)
\begin{itemize}[noitemsep]
  \item Un nod transmite; dacă este primit, orice nod funcțional retransmite
  \item Toate nodurile funcționale transmit același mesaj în aceeași ordine
\end{itemize}
\end{formulabox}

\subsection{Dificultăți în procesarea paralelă}
\begin{enumerate}[noitemsep]
  \item \textbf{Debugging} — greu de replicat contextul (cauza: \emph{interleaving})
  \item \textbf{Corectitudinea} — demonstrarea formală este dificilă
  \item \textbf{Benchmark} — evaluarea pe input mic poate induce în eroare
  \item \textbf{Comunicarea} — proiectarea optimă inter-procese
\end{enumerate}

% =============================================================================
\section{Curs 2 — Tipuri și Nivele de Paralelism + POSIX}
% =============================================================================

\subsection{Taxonomia lui Flynn (1972)}

\begin{center}
\begin{tabular}{ccll}
\toprule
\thead{Tip} & \thead{Instrucțiuni} & \thead{Date} & \thead{Descriere} \\
\midrule
\textbf{SISD} & 1 flux & 1 flux & Calculator clasic secvențial \\
\textbf{SIMD} & 1 flux & Multiple fluxuri & GPU, procesare vectorială \\
\textbf{MISD} & Multiple fluxuri & 1 flux & Teoretic, rar folosit \\
\textbf{MIMD} & Multiple fluxuri & Multiple fluxuri & Cel mai comun: sisteme multiprocesor \\
\bottomrule
\end{tabular}
\end{center}

\begin{itemize}[noitemsep]
  \item \textbf{SIMD}: O singură unitate de control, procesoare diferite primesc aceleași instrucțiuni pe date diferite. Rigid, dar eficient pentru date uniform structurate.
  \item \textbf{MIMD}: Cel mai flexibil. Fiecare procesor are flux propriu de instrucțiuni și date proprii.
\end{itemize}

\subsection{PRAM (Parallel Random Access Machine)}

Model teoretic cu memorie partajată unde accesul se face în \textbf{timp constant} (cost 0). Omologul modelului von Neumann pentru sisteme paralele.

\textbf{Instrucțiune PRAM}: \lstinline|For all x in X do in parallel #instruction(x)|

\subsubsection{Variante PRAM (după accesul la memorie):}

\begin{center}
\begin{tabular}{lll}
\toprule
\thead{Variantă} & \thead{Citire} & \thead{Scriere} \\
\midrule
\textbf{EREW} & Exclusivă & Exclusivă \\
\textbf{CREW} & Concurentă & Exclusivă \\
\textbf{ERCW} & Exclusivă & Concurentă \\
\textbf{CRCW} & Concurentă & Concurentă \\
\bottomrule
\end{tabular}
\end{center}

EREW = cel mai realist \quad $\longrightarrow$ \quad CRCW = cel mai puțin realist

\subsection{Metrici rețele de interconectare}

\begin{center}
\begin{tabular}{lcp{7cm}}
\toprule
\thead{Metrică} & \thead{Direcție} & \thead{Definiție} \\
\midrule
\textbf{Diametru} & $\downarrow$ & Lungimea maximă dintre oricare 2 noduri \\
\textbf{Conectivitate} & $\uparrow$ & Nr. minim de conexiuni de distrus pt. 2 componente \\
\textbf{Lăț. bisecțiunii} & $\uparrow$ & Nr. minim de muchii ce taie rețeaua în 2 jumătăți egale \\
\textbf{Cost} & $\downarrow$ & Numărul total de legături \\
\bottomrule
\end{tabular}
\end{center}

\begin{center}
\begin{tabular}{lcccc}
\toprule
\thead{Topologie} & \thead{Diametru} & \thead{Conectivitate} & \thead{Lăț. bisecțiunii} & \thead{Cost} \\
\midrule
Bus & $O(1)$ & $O(1)$ & $O(1)$ & $O(p)$ \\
Crossbar & $O(1)$ & $O(1)$ & $O(p)$ & $O(p^2)$ \\
Hipercub & $O(\log p)$ & $O(\log p)$ & $O(p/2)$ & $O(p \log p)$ \\
\bottomrule
\end{tabular}
\end{center}

\subsection{Costul comunicării}

\begin{formulabox}
\[
T_{com} = T_s + T_w \cdot m
\]
\begin{itemize}[noitemsep]
  \item $T_s$ = timpul de inițializare (setup, routing)
  \item $T_w$ = timpul de transfer per cuvânt $= 1 / \text{lățimea canalului}$
  \item $m$ = dimensiunea mesajului (în cuvinte)
\end{itemize}
\textbf{Important:} $T_s \gg T_w$ în practică $\Rightarrow$ reduceți numărul de mesaje, nu dimensiunea lor!
\end{formulabox}

\subsection{Accelerarea (Speedup)}

\begin{formulabox}
\[
S(p) = \frac{T_1}{T_p}
\]
$T_1$ = cel mai bun timp secvențial pe 1 procesor; \; $T_p$ = timp paralel pe $p$ procesoare.

\textbf{Atenție}: Se folosește \emph{cel mai bun} algoritm secvențial!
\end{formulabox}

\subsection{Pipelining}

\begin{itemize}[noitemsep]
  \item Se descompune procesul în $K$ \textbf{stadii/segmente} distincte
  \item Fiecare stadiu este dedicat unui procesor
  \item Limitare: cel mai lent stadiu determină ritmul întregului pipeline
\end{itemize}

\begin{formulabox}
\[
T_k = (K + N - 1) \cdot T \qquad S = \frac{N \cdot K}{K + N - 1}
\]
\begin{itemize}[noitemsep]
  \item $K$ = numărul de stadii
  \item $N$ = numărul de elemente de procesat
  \item $T$ = durata unui ciclu (= cel mai lent stadiu)
  \item Pentru $N \gg K$: $S \approx K$
\end{itemize}
\end{formulabox}

% =============================================================================
\section{Curs 3 — Modele de Calcul Paralel + C++ Threads}
% =============================================================================

\subsection{Paralelism implicit vs. explicit}

\begin{itemize}[noitemsep]
  \item \textbf{Implicit}: compilatorul paralelizează automat
  \item \textbf{Explicit}: programatorul specifică paralelismul (directive, biblioteci)
\end{itemize}

\begin{center}
\begin{tabular}{lc}
\toprule
\thead{Instrucțiuni} & \thead{Pot fi paralele?} \\
\midrule
\lstinline|a = 2; b = 3;| & \textcolor{codegreen}{\textbf{DA}} — independente \\
\lstinline|a = 3; b = a;| & \textcolor{accent}{\textbf{NU}} — dependență de date \\
\lstinline|for (i=0;i<100;i++) b[i]=i;| & \textcolor{codegreen}{\textbf{DA}} — iterații independente \\
\lstinline|a = f(x); b = g(x);| & \textcolor{accent}{\textbf{Depinde}} de funcțiile $f$ și $g$ \\
\bottomrule
\end{tabular}
\end{center}

\subsection{Graful de dependințe (DAG)}

\begin{itemize}[noitemsep]
  \item \textbf{DAG} = Graf Direcționat Aciclic
  \item Noduri = task-uri; Muchii = dependențe ($A \to B$ înseamnă B depinde de A)
  \item \textbf{Calea critică} = cel mai lung drum în DAG $\Rightarrow$ limita inferioară a timpului paralel
\end{itemize}

\subsection{Metode de descompunere}

\begin{center}
\begin{tabularx}{\textwidth}{lX}
\toprule
\thead{Metodă} & \thead{Când se folosește} \\
\midrule
\textbf{Descompunerea datelor} & Input/output voluminos; partiționăm datele $\to$ task-uri \\
\textbf{Descompunerea task-urilor} & Procesări distincte care pot rula independent \\
\textbf{Recursivă} & Divide et impera: quicksort, min, max \\
\textbf{Exploratorie} & Spațiu de soluții; poate produce \emph{speedup superliniar} \\
\textbf{Speculativă} & Pașii următori depind de rezultatul curent; se execută predicții \\
\textbf{Hibridă} & Combinație a metodelor de mai sus \\
\bottomrule
\end{tabularx}
\end{center}

\subsection{Modelele algoritmice}

\begin{center}
\begin{tabularx}{\textwidth}{lX}
\toprule
\thead{Model} & \thead{Principiu} \\
\midrule
\textbf{Paralelizarea datelor} & Input împărțit între procesoare; segmente diferite procesate în paralel \\
\textbf{Paralelizare cu grafuri} & Algoritmul descompus în secțiuni distincte (fork/join/spawn) \\
\textbf{Paralelizare cu pool} & Task-urile asignate dinamic; fiecare task are puține date \\
\bottomrule
\end{tabularx}
\end{center}

\subsection{Alocarea task-urilor}

\begin{formulabox}
\[
T_o = p \cdot T_p - T_1 \quad \text{(overhead total)}
\]
\end{formulabox}

\textbf{Surse de overhead}: dezechilibrul încărcării, comunicarea, sincronizarea.

\textbf{Strategii (pot fi conflictuale):}
\begin{itemize}[noitemsep]
  \item Task-uri independente $\to$ procesoare \textbf{diferite} ($\uparrow$ concurență)
  \item Task-uri care comunică frecvent $\to$ \textbf{același} procesor ($\uparrow$ localizare)
\end{itemize}

% =============================================================================
\section{Curs 4 — Metrici de Performanță + STL Parallel}
% =============================================================================

\subsection{Metrici fundamentale}

\begin{formulabox}
\begin{align*}
\text{Accelerare:} \quad & S(p) = \frac{T_1}{T_p} \\[4pt]
\text{Eficiență:} \quad & E(p) = \frac{S(p)}{p} = \frac{T_1}{p \cdot T_p} \quad (0 < E \leq 1) \\[4pt]
\text{Cost total:} \quad & C = p \cdot T_p \\[4pt]
\text{Overhead:} \quad & T_o = p \cdot T_p - T_1
\end{align*}
\end{formulabox}

\subsection{Legea lui Amdahl}

\begin{defbox}{Legea lui Amdahl — dimensiunea problemei este FIXĂ}
\[
S(N) = \frac{1}{\sigma + \dfrac{1 - \sigma}{N}}
\]
\begin{itemize}[noitemsep]
  \item $\sigma$ = fracția din cod care \emph{trebuie} executată secvențial \quad ($0 < \sigma \leq 1$)
  \item $N$ = numărul de procesoare
  \item $S(\infty) = \dfrac{1}{\sigma}$ — \textbf{limita maximă absolută!}
\end{itemize}
\end{defbox}

\begin{center}
\begin{tabular}{ccc}
\toprule
\thead{$\sigma$ (cod serial)} & \thead{$S(\infty)$} & \thead{$S(100)$} \\
\midrule
10\% & 10 & 9.17 \\
30\% & 3.33 & 3.23 \\
1\% & 100 & 50.25 \\
\bottomrule
\end{tabular}
\end{center}

\begin{tipbox}
Chiar și \textbf{10\% cod serial} limitează accelerarea la maxim \textbf{10×}, indiferent de câte procesoare avem!

\textbf{Limitări}: Nu ia în calcul overhead-ul de comunicare/sincronizare și nici dimensiunea problemei.
\end{tipbox}

\subsection{Legea lui Gustafson (Accelerarea Scalată)}

\begin{defbox}{Legea lui Gustafson — dimensiunea problemei CREȘTE cu numărul de procesoare}
\[
S \leq p + (1 - p) \cdot s
\]
\begin{itemize}[noitemsep]
  \item $p$ = numărul de procesoare
  \item $s$ = procentul de execuție serială \emph{din aplicația paralelă}
\end{itemize}
\end{defbox}

\textbf{Exemplu}: 32 procesoare, 1\% serial:
\begin{itemize}[noitemsep]
  \item Gustafson: $S \leq 32 + (1 - 32) \cdot 0.01 = 31.69$
  \item Amdahl: $S \leq 24.43$ (mai pesimist și nerealist pentru probleme mari)
\end{itemize}

\subsection{STL Parallel (C++17)}

69 de algoritmi paraleli în biblioteca standard. Paralelizarea = adăugarea unui \textbf{singur parametru}:

\begin{lstlisting}[language=C++]
#include <execution>
std::sort(std::execution::par, begin(v), end(v));  // paralel
\end{lstlisting}

\begin{center}
\begin{tabular}{ll}
\toprule
\thead{Politică} & \thead{Semnificație} \\
\midrule
\lstinline|std::execution::seq| & Secvențial \\
\lstinline|std::execution::par| & Paralel (fără vectorizare) \\
\lstinline|std::execution::par_unseq| & Paralel + vectorizat \\
\lstinline|std::execution::unseq| & Vectorizat fără paralelism (C++20) \\
\bottomrule
\end{tabular}
\end{center}

\begin{tipbox}
Politica este un \emph{hint} — biblioteca poate ignora și executa serial.
\textbf{NU} rezolvă probleme de concurență globale! Evitați \emph{data race}-urile explicit.
\end{tipbox}

% =============================================================================
\section{Curs 5 — Arhitecturi de Calcul Paralel/Distribuit + Java Paralel}
% =============================================================================

\subsection{Arhitecturi cu memorie partajată}

\subsubsection{UMA — Uniform Memory Access (SMP)}
\begin{itemize}[noitemsep]
  \item Toate procesoarele accesează memoria cu \textbf{același timp}
  \item \textbf{CC-UMA} = Cache Coherent UMA (coerența cache asigurată HW)
  \item Limitat la câteva zeci de procesoare
  \item Exemple: Intel Core i7, mașini SMP clasice
\end{itemize}

\subsubsection{NUMA — Non-Uniform Memory Access}
\begin{itemize}[noitemsep]
  \item Obținut prin conectarea mai multor SMP-uri
  \item Accesul la memoria altui SMP are timp \textbf{mai mare}
  \item \textbf{CC-NUMA} = Cache Coherent NUMA
  \item Mai scalabil decât UMA
\end{itemize}

\subsubsection{Comparație arhitecturi}

\begin{center}
\begin{tabular}{lccc}
\toprule
\thead{} & \thead{CC-UMA} & \thead{CC-NUMA} & \thead{Distribuită} \\
\midrule
Comunicare & MPI, Threads & MPI, OpenMP & MPI \\
Scalabilitate & $\times 10$ proc. & $\times 100$ proc. & $\times 1000$ proc. \\
Exemple & SMP & SGI Origin & IBM SP2, BlueGene \\
\bottomrule
\end{tabular}
\end{center}

\begin{itemize}[noitemsep]
  \item \textbf{Memorie partajată} — Avantaj: ușor de programat. Dezavantaj: scalabilitate limitată.
  \item \textbf{Memorie distribuită} — Avantaj: scalabilitate excelentă. Dezavantaj: programatorul gestionează explicit comunicarea.
  \item \textbf{Hibrid} (cluster de SMP-uri) — tendința actuală în HPC.
\end{itemize}

\subsection{Arhitecturi sisteme distribuite}

\subsubsection{Client-Server}
\begin{itemize}[noitemsep]
  \item Clienții cunosc serverele; serverele \textbf{NU} cunosc clienții
  \item \textbf{3 nivele}: View (client) $\to$ Controller (server aplicație) $\to$ Model (baza de date)
  \item Clase: Host-based, Server-based, Client-based, \textbf{Cooperative}
\end{itemize}

\subsubsection{Fiabilitatea sistemelor distribuite}
\begin{center}
\begin{tabular}{lll}
\toprule
\thead{} & \thead{Cu blocare (sincron)} & \thead{Fără blocare (asincron)} \\
\midrule
Comportament & Procesul așteaptă confirmarea & Procesul continuă imediat \\
Avantaje & Consistență garantată & Eficient și flexibil \\
Dezavantaje & Poate bloca mult & Debugging dificil \\
\bottomrule
\end{tabular}
\end{center}

\subsection{Java Paralel — Fork/Join Framework}

\begin{center}
\begin{tabularx}{\textwidth}{lX}
\toprule
\thead{Clasă} & \thead{Rol} \\
\midrule
\lstinline|ForkJoinTask<V>| & Clasă abstractă de bază \\
\lstinline|ForkJoinPool| & Administrează și monitorizează execuția task-urilor \\
\lstinline|RecursiveAction| & Task fără rezultat (void) \\
\lstinline|RecursiveTask<V>| & Task cu rezultat \\
\bottomrule
\end{tabularx}
\end{center}

\textbf{Strategie}: Divide and Conquer recursiv — împarte până când task-ul e suficient de mic, apoi execuție secvențială.

\begin{lstlisting}[language=Java]
class MyTask extends RecursiveAction {
    protected void compute() {
        if (dimensiuneMica) {
            // executa secvential
        } else {
            MyTask left = new MyTask(...);
            MyTask right = new MyTask(...);
            left.fork();   // executie asincrona
            right.fork();
            left.join();   // asteapta terminarea
            right.join();
        }
    }
}
\end{lstlisting}

\subsection{Programare concurentă Java}

\begin{lstlisting}[language=Java]
// synchronized: un singur fir executa metoda la un moment dat
public synchronized void metoda() { ... }

// ExecutorService (nivel inalt)
ExecutorService exec = Executors.newFixedThreadPool(4);
exec.execute(new RunnableTask());
exec.shutdown();

// Fluxuri paralele
list.parallelStream()
    .filter(x -> x > 0)
    .collect(Collectors.toList());
\end{lstlisting}

% =============================================================================
\section{Curs 6 — Proiectarea Programelor Paralele + OpenMP}
% =============================================================================

\subsection{Metodologia Foster — PCAM}

\begin{defbox}{Cei 4 pași Foster}
\begin{enumerate}[noitemsep]
  \item \textbf{Partiționare (P)} — granularitate fină, maximizăm potențialul de paralelism
  \item \textbf{Comunicare (C)} — identificăm fluxurile de date între task-uri
  \item \textbf{Aglomerare (A)} — grupăm task-uri mici în task-uri mai mari
  \item \textbf{Alocare/Mapare (M)} — asignăm task-uri pe procesoare fizice
\end{enumerate}
\end{defbox}

\subsection{Pasul 1: Partiționarea (Descompunerea)}

\begin{itemize}[noitemsep]
  \item \textbf{Descompunerea domeniului} (a datelor): datele $\to$ task-uri
  \item \textbf{Descompunerea funcțională}: funcțiile $\to$ task-uri
  \item Sunt \textbf{complementare}; se pot aplica la aceeași problemă
\end{itemize}

\textbf{Verificare corectitudine:}
\begin{enumerate}[noitemsep]
  \item Nr. task-uri $\gg$ nr. procesoare? (cel puțin $10\times$)
  \item Evită procesări/date redundante?
  \item Task-urile au dimensiuni comparabile?
  \item Nr. task-uri scalează cu dimensiunea problemei?
\end{enumerate}

\subsection{Pasul 2: Comunicarea}

\begin{center}
\begin{tabular}{ll}
\toprule
\thead{Dimensiune} & \thead{Tipuri} \\
\midrule
Scope & Local (vecini imediați) vs. Global (toate task-urile) \\
Structură & Structurată (grid/arbore) vs. Nestructurată (graf oarecare) \\
Timing & Statică (parteneri ficși) vs. Dinamică (schimbătoare la runtime) \\
Coordonare & Sincronă (coordonat) vs. Asincronă (independent) \\
\bottomrule
\end{tabular}
\end{center}

\subsubsection{Exemplu: Suma N numere — Divide and Conquer}

Naive (serial): $O(N)$ — \textbf{rău}. \quad Divide and Conquer: $O(\log N)$ — \textbf{bun}!

\[
\text{Nivel 0: } [a_0][a_1][a_2][a_3]\cdots \xrightarrow{\text{nivel 1}} [a_0+a_1][a_2+a_3]\cdots \xrightarrow{\text{nivel 2}} \cdots \xrightarrow{\log N \text{ pași}} \text{Sum}
\]

La același nivel toate adunările se execută \textbf{concurent}!

\subsection{Pasul 3: Aglomerarea}

\begin{itemize}[noitemsep]
  \item Grupăm task-uri mici $\to$ task-uri mai mari ($\downarrow$ comunicare, $\uparrow$ granularitate)
  \item \textbf{Structura fluture (butterfly)}: $\log N$ nivele; suma completă în $\log N$ pași
\end{itemize}

\subsection{Pasul 4: Alocarea}

\begin{tipbox}
Problema alocării este \textbf{NP-completă} în cazul general!
\end{tipbox}

\begin{itemize}[noitemsep]
  \item Strategii conflictuale:
    \begin{itemize}[noitemsep]
      \item Task-uri independente $\to$ procesoare \textbf{diferite} ($\uparrow$ concurență)
      \item Task-uri care comunică frecvent $\to$ \textbf{același} procesor ($\uparrow$ localizare)
    \end{itemize}
\end{itemize}

\subsubsection{Algoritmi de echilibrare a încărcării}

\begin{center}
\begin{tabularx}{\textwidth}{lX}
\toprule
\thead{Algoritm} & \thead{Descriere} \\
\midrule
\textbf{Bisecțiunea recursivă} & Divide domeniul în subdomenii de cost aproximativ egal \\
\textbf{Local (Greedy)} & Compară cu vecinii; transferă task-uri dacă diferența $>$ prag \\
\textbf{Probabilistic} & Alocă aleator; bun când $N_{\text{task}} \gg N_{\text{proc}}$ \\
\textbf{Manager/Worker} & Manager distribuie task-uri; Worker-ii execută și cer mai mult \\
\bottomrule
\end{tabularx}
\end{center}

\subsection{Concluzii Metodologie Foster}

\begin{center}
\begin{tabular}{lcl}
\toprule
\thead{Pas} & \thead{Depinde de arhitectură?} & \thead{Scop} \\
\midrule
Descompunere & NU & Permite concurența \\
Aglomerare & NU & Echilibrează încărcarea, limitează comunicarea \\
Alocare & \textbf{DA} & Exploatează localizarea, reduce comunicarea \\
\bottomrule
\end{tabular}
\end{center}

\subsection{OpenMP}

\subsubsection{Ce este OpenMP?}
Set de \textbf{directive compilator} + funcții de bibliotecă pentru C/C++/Fortran. Model: \textbf{Fork-Join} cu memorie partajată. Paralelism adăugat \textbf{incremental} la cod secvențial.

\subsubsection{Sintaxa de bază}

\begin{lstlisting}[language=C++]
#include <omp.h>

#pragma omp parallel num_threads(4)
{
    int id = omp_get_thread_num();    // ID-ul firului curent
    int n  = omp_get_num_threads();   // numarul total de fire
    printf("Thread %d din %d\n", id, n);
}
\end{lstlisting}

\subsubsection{Sincronizare în OpenMP}

\begin{lstlisting}[language=C++]
// 1. Barrier: toate firele asteapta
#pragma omp barrier

// 2. Critical: excludere mutuala (sectiune critica)
#pragma omp critical
{
    result += local_value;   // un singur fir la un moment dat
}

// 3. Atomic: update atomic al unei singure variabile (mai rapid)
#pragma omp atomic
X += tmp;
\end{lstlisting}

\subsubsection{Paralelizarea buclelor}

\begin{lstlisting}[language=C++]
#pragma omp parallel
{
    #pragma omp for
    for (int i = 0; i < N; i++) {
        A[i] = f(i);   // variabila i este implicit "privata"
    }
}
\end{lstlisting}

\subsubsection{Opțiuni de planificare (schedule)}

\begin{center}
\begin{tabularx}{\textwidth}{lX}
\toprule
\thead{Schedule} & \thead{Comportament} \\
\midrule
\lstinline|static [,chunk]| & Blocuri de dimensiune \textit{chunk} distribuite round-robin \\
\lstinline|dynamic [,chunk]| & Fiecare fir ia dinamic \textit{chunk} iterații din coadă \\
\lstinline|guided [,chunk]| & Blocuri descrescătoare, luate dinamic \\
\lstinline|runtime| & Planificarea din variabila de mediu \texttt{OMP\_SCHEDULE} \\
\lstinline|auto| & La dispoziția compilatorului/runtime-ului \\
\bottomrule
\end{tabularx}
\end{center}

\begin{tipbox}
\textbf{Data race} în OpenMP: apar când mai multe fire accesează aceeași variabilă și cel puțin unul scrie.
Soluții: \lstinline|critical|, \lstinline|atomic|, variabile private per fir. \textbf{Sincronizarea este costisitoare} — minimizați-o!
\end{tipbox}

% =============================================================================
\section{Rezumat Formule Esențiale}
% =============================================================================

\begin{formulabox}
\begin{align*}
\text{Accelerare:} \quad & S(p) = \frac{T_1}{T_p} \\[6pt]
\text{Eficiență:} \quad & E(p) = \frac{S(p)}{p} \quad (0 < E \leq 1) \\[6pt]
\text{Cost comunicare:} \quad & T_{com} = T_s + T_w \cdot m \\[6pt]
\text{Timp pipeline:} \quad & T_k = (K + N - 1) \cdot T \\[6pt]
\text{Amdahl:} \quad & S_A = \frac{1}{\sigma + \dfrac{1-\sigma}{N}}, \quad S(\infty) = \frac{1}{\sigma} \\[8pt]
\text{Gustafson:} \quad & S_G \leq p + (1-p) \cdot s \\[6pt]
\text{Overhead:} \quad & T_o = p \cdot T_p - T_1
\end{align*}
\end{formulabox}

% =============================================================================
\section{Tips pentru Examen}
% =============================================================================

\begin{enumerate}
  \item \textbf{Taxonomia Flynn}: știți toate 4 tipuri cu exemple clare (SIMD = GPU/vectorial, MIMD = multiprocesor).
  \item \textbf{Legea Amdahl}: calculați $S(\infty) = 1/\sigma$ și pentru valori concrete de $N$.
  \item \textbf{Metodologia Foster}: cei 4 pași \textbf{în ordine} și ce face fiecare.
  \item \textbf{UMA vs. NUMA vs. Distribuită}: scalabilitate $\times 10$, $\times 100$, $\times 1000$.
  \item \textbf{Concurență $\neq$ Paralelism}: definiție exactă pentru ambele.
  \item \textbf{Data race}: ce este, cum apare, cum se previne (critical/atomic/barrier).
  \item \textbf{OpenMP}: directivele principale și când se folosesc.
  \item \textbf{Problema celor 2 generali}: nu are soluție; transmiterea atomică $\equiv$ consensus.
  \item \textbf{Calea critică}: limitează accelerarea maximă obtenabilă (limita inferioară a timpului paralel).
  \item \textbf{Amdahl vs. Gustafson}: Amdahl = pesimist pentru dimensiune fixă; Gustafson = realist când problema crește.
\end{enumerate}

\end{document}
